\documentclass[11pt, letterpaper]{article}

% ==============================================================================
% PACKAGES
% ==============================================================================
\usepackage{amsmath, amssymb, amsthm}
\usepackage{graphicx}
\usepackage{hyperref}
\usepackage{geometry}
\usepackage{listings}
\usepackage{booktabs}
\usepackage{enumitem}
\usepackage{xcolor}
\usepackage{algorithm}
\usepackage{algorithmic}

\geometry{margin=1in}

% Hyperlink styling
\hypersetup{
    colorlinks=true,
    linkcolor=blue,
    filecolor=magenta,
    urlcolor=cyan,
    citecolor=blue,
}

% Code listing styling
\lstset{
    basicstyle=\ttfamily\small,
    breaklines=true,
    columns=fullflexible,
    frame=single,
    backgroundcolor=\color{gray!10},
}

% Theorem environments
\newtheorem{axiom}{Axiom}
\newtheorem{theorem}{Theorem}
\newtheorem{lemma}{Lemma}
\newtheorem{definition}{Definition}
\newtheorem{property}{Property}

% ==============================================================================
% METADATA
% ==============================================================================
\title{\textbf{ORACLE SUPERTEAM}:\\
A Governance-Centric Multi-Agent Framework\\
for Deterministic, Evidence-Bound Decision-Making}

\author{
Jean Marie Tassy Simeoni\\
\small JMT Consulting\\
\small \texttt{contact@jmtconsulting.com}
}

\date{\today\\[1em]
\small \textit{Version 1.0 — Constitutional Model}}

% ==============================================================================
% DOCUMENT
% ==============================================================================
\begin{document}

\maketitle

% ==============================================================================
% ABSTRACT
% ==============================================================================
\begin{abstract}
\noindent
This paper introduces \textbf{ORACLE SUPERTEAM}, a governance-first multi-agent framework designed to transform loosely structured human or artificial reasoning into auditable, deterministic, and institutionally reliable decisions. Unlike conventional multi-agent systems that emphasize emergent consensus, narrative collaboration, or optimization through averaging, ORACLE SUPERTEAM is grounded in \textit{epistemic sovereignty}: no claim may be accepted without explicit, verifiable evidence, and no agent—human or artificial—possesses unilateral authority to decide outcomes.

The framework formalizes decision-making as a jurisdictional pipeline composed of three strictly separated layers: \textbf{(i) production}, where specialized agent teams (``Superteams'') generate candidate obligations and non-binding signals; \textbf{(ii) adjudication}, where a critic agent applies lexicographic veto rules based on evidence sufficiency, logical coherence, and real-world impact; and \textbf{(iii) integration}, where a deterministic automaton emits a binary verdict (\texttt{SHIP} / \texttt{NO\_SHIP}).

Central to the system is the principle \texttt{NO\_RECEIPT = NO\_SHIP}, ensuring that any claim asserting real-world effects or attestable facts is blocked unless accompanied by cryptographically hashable proof.

ORACLE SUPERTEAM departs from mainstream agentic architectures by rejecting conversational consensus, probabilistic voting, and opaque ``reasoning traces'' as decision criteria. Instead, it relies on event-sourced memory, deterministic hashing, and explicit obligation semantics to guarantee replayability, auditability, and resistance to rhetorical inflation.

We present the complete implementation, formal semantics, threat model, and test vault validation results. The framework demonstrates that scalable collective intelligence does not emerge from more dialogue, but from \textit{hard constraints, explicit evidence, and irreversible records}—reframing multi-agent systems as institutions rather than conversations.
\end{abstract}

% ==============================================================================
% 1. INTRODUCTION
% ==============================================================================
\section{Introduction}

Multi-agent systems have emerged as a promising paradigm for tackling complex decision-making tasks requiring diverse expertise \cite{autogen, chatdev, metagpt}. However, most current architectures suffer from fundamental limitations:

\begin{itemize}[leftmargin=*]
    \item \textbf{Non-determinism}: LLM-based agents produce variable outputs for identical inputs, making replay and audit impossible.
    \item \textbf{Sovereignty leaks}: Individual agents can influence decisions through persuasive language rather than verifiable evidence.
    \item \textbf{Consensus drift}: Averaging or voting mechanisms can be gamed through rhetorical inflation or strategic voting.
    \item \textbf{Lack of hard constraints}: Most systems have no enforced veto mechanisms for safety-critical concerns.
\end{itemize}

ORACLE SUPERTEAM addresses these limitations by treating multi-agent decision-making not as a conversation, but as a \textit{computational tribunal}—a constitutional system with explicit roles, deterministic rules, and cryptographic auditability.

\subsection{Motivating Principles}

The design is guided by three core insights:

\begin{enumerate}
    \item \textbf{Evidence over rhetoric}: Claims must be supported by attestable, hashable proof (``receipts''), not persuasive arguments.
    \item \textbf{Separation of production and adjudication}: Agent teams produce signals and obligations but do not vote or decide outcomes.
    \item \textbf{Binary verdicts with lexicographic veto}: Decisions resolve to exactly two states (\texttt{SHIP} / \texttt{NO\_SHIP}), with kill-switches overriding all other signals.
\end{enumerate}

\subsection{Contributions}

This paper makes the following contributions:

\begin{enumerate}
    \item A formal model for multi-agent governance based on jurisdictional separation and deterministic adjudication.
    \item Implementation of a complete, zero-dependency Python framework with 100\% test vault coverage.
    \item Demonstration of \textit{replay determinism}—identical inputs producing identical hashes across all executions.
    \item A threat model addressing sovereignty leaks, consensus gaming, and audit failure modes.
    \item Open-source release with constitutional documentation and visualization tooling.
\end{enumerate}

% ==============================================================================
% 2. SYSTEM AXIOMS
% ==============================================================================
\section{Constitutional Axioms}

ORACLE SUPERTEAM enforces five non-negotiable axioms:

\begin{axiom}[\texttt{NO\_RECEIPT = NO\_SHIP}]
\label{ax:no-receipt}
All Tier I claims (asserting real-world effects) require cryptographically hashable evidence. Claims without receipts are automatically blocked.
\end{axiom}

\begin{axiom}[Non-Sovereign Production]
\label{ax:non-sovereign}
Superteams (agent teams) may emit signals or obligations, but they may \textit{never} vote, score, or directly influence verdicts. All decision authority resides in the adjudication and integration layers.
\end{axiom}

\begin{axiom}[Binary Verdict Semantics]
\label{ax:binary}
The system outputs exactly two verdict states: \texttt{SHIP} (ship\_permitted = true) and \texttt{NO\_SHIP} (ship\_permitted = false). Internal classifications (ACCEPT, QUARANTINE, KILL) are diagnostic only.
\end{axiom}

\begin{axiom}[Lexicographic Veto Dominance]
\label{ax:kill-switch}
Any authorized \texttt{KILL\_REQUEST} signal from designated teams (Legal Office, Security Sector) overrides all other signals, scores, and evaluations. Kill-switches are immediate and non-negotiable.
\end{axiom}

\begin{axiom}[Replay Determinism]
\label{ax:replay}
Identical inputs (claim + evidence + signals + configuration) must produce identical outputs (obligations + verdict + hashes) across all executions. Non-determinism indicates system corruption.
\end{axiom}

% ==============================================================================
% 3. FORMAL MODEL
% ==============================================================================
\section{Formal Model}

\subsection{Claims}

\begin{definition}[Claim]
A claim $C$ is a 6-tuple:
\[
C = \langle id, \text{assertion}, \text{tier}, \text{domain}, \text{requires\_receipts}, \text{owner\_team} \rangle
\]
where:
\begin{itemize}
    \item $id \in \mathbb{S}$ is a unique identifier
    \item $\text{assertion} \in \mathbb{S}$ is the statement being evaluated
    \item $\text{tier} \in \{\text{I}, \text{II}, \text{III}\}$ indicates impact level
    \item $\text{domain} \subseteq \mathcal{D}$ is a set of domain tags
    \item $\text{requires\_receipts} \in \{\text{true}, \text{false}\}$
    \item $\text{owner\_team} \in \mathcal{T}$ is the submitting team
\end{itemize}
\end{definition}

Claims are \textit{immutable}: once submitted, they cannot be modified. All changes require new claim submissions.

\subsection{Signals}

\begin{definition}[Signal]
A signal $s$ is a 4-tuple emitted by an agent team:
\[
s = \langle \text{team}, \text{signal\_type}, \text{obligation\_type}, \text{rationale} \rangle
\]
where:
\begin{itemize}
    \item $\text{team} \in \mathcal{T}$ is the emitting team
    \item $\text{signal\_type} \in \{\text{OBLIGATION\_OPEN}, \text{RISK\_FLAG}, \text{EVIDENCE\_WEAK}, \text{KILL\_REQUEST}\}$
    \item $\text{obligation\_type} \in \mathcal{O} \cup \{\bot\}$ (optional, required for \texttt{OBLIGATION\_OPEN})
    \item $\text{rationale} \in \mathbb{S}$ is human-readable justification
\end{itemize}
\end{definition}

\textbf{Critical property}: Signals are \textit{non-sovereign}. They request obligations or flag concerns but do not directly influence verdicts.

\subsection{Obligations}

\begin{definition}[Obligation]
An obligation $o$ is a blocking constraint:
\[
o = \langle \text{type}, \text{owner\_team}, \text{closure\_criteria}, \text{status} \rangle
\]
where:
\begin{itemize}
    \item $\text{type} \in \{\text{METRICS\_INSUFFICIENT}, \text{LEGAL\_COMPLIANCE}, \text{SECURITY\_THREAT}, \text{EVIDENCE\_MISSING}, \text{CONTRADICTION\_DETECTED}\}$
    \item $\text{owner\_team} \in \mathcal{T}$ is responsible for resolution
    \item $\text{closure\_criteria} \in \mathbb{S}$ specifies what evidence clears the obligation
    \item $\text{status} \in \{\text{OPEN}, \text{RESOLVED}, \text{INVALIDATED}\}$
\end{itemize}
\end{definition}

\begin{property}[Obligation Blocking]
If $\exists o \in \mathcal{O}_{\text{open}}$ such that $o.\text{status} = \text{OPEN}$, then the verdict must be \texttt{NO\_SHIP}.
\end{property}

\subsection{Evidence and Receipts}

\begin{definition}[Evidence]
Evidence $e$ is attestable proof:
\[
e = \langle id, \text{type}, \text{tags}, \text{hash}, \text{issuer} \rangle
\]
where:
\begin{itemize}
    \item $\text{hash} \in \mathbb{H}$ is a SHA-256 digest (when applicable)
    \item $\text{tags} \subseteq \mathcal{TAG}$ are semantic markers for contradiction detection
\end{itemize}
\end{definition}

\begin{definition}[Receipt]
A receipt $r$ is evidence that can be cryptographically verified:
\[
\text{verify}(r) = (\text{hash}(r.\text{content}) = r.\text{claimed\_hash})
\]
\end{definition}

% ==============================================================================
% 4. ARCHITECTURAL LAYERS
% ==============================================================================
\section{Three-Layer Architecture}

\subsection{Layer 1: Production (Non-Sovereign)}

The production layer consists of specialized agent teams (Superteams) organized by domain:

\begin{itemize}
    \item \textbf{Engineering Wing}: Technical feasibility, metrics validation
    \item \textbf{Legal Office}: Compliance, regulatory requirements (kill-switch authority)
    \item \textbf{Security Sector}: Threat assessment, vulnerability analysis (kill-switch authority)
    \item \textbf{Data Validation Office}: Evidence quality, statistical rigor
    \item \textbf{Strategy HQ}: Business alignment, strategic fit
    \item \textbf{COO HQ}: Operational readiness, resource availability
    \item \textbf{UX/Impact Bureau}: User impact, accessibility
\end{itemize}

Each team emits \textit{signals}, not votes:

\begin{lstlisting}[language=Python, caption=Signal emission example]
{
  "team": "Engineering Wing",
  "signal_type": "OBLIGATION_OPEN",
  "obligation_type": "METRICS_INSUFFICIENT",
  "rationale": "Need A/B test results with statistical significance."
}
\end{lstlisting}

\subsection{Layer 2: Adjudication (Lexicographic Veto)}

The adjudication layer applies deterministic rules in strict priority:

\begin{algorithm}
\caption{Adjudication Pipeline}
\begin{algorithmic}[1]
\STATE \textbf{Input:} Claim $C$, Evidence $E$, Signals $S$
\STATE \textbf{Output:} Adjudication state $A$
\STATE
\STATE $\text{kill\_triggered} \gets \text{CheckKillSignals}(S)$
\STATE $\text{contradictions} \gets \text{DetectContradictions}(E)$
\STATE $\text{rule\_kill} \gets \text{CheckHighSeverity}(\text{contradictions})$
\STATE $\text{obligations} \gets \text{GenerateObligations}(S, \text{contradictions})$
\STATE
\RETURN $A = \langle \text{kill\_triggered}, \text{rule\_kill}, \text{contradictions}, \text{obligations} \rangle$
\end{algorithmic}
\end{algorithm}

\textbf{Contradiction Detection}: Hard-coded rules check evidence tags for logical conflicts:

\begin{itemize}
    \item \textbf{HC-PRIV-001} (HIGH): ``no\_personal\_data\_claim'' + ``biometric'' $\Rightarrow$ AUTO-KILL
    \item \textbf{HC-SEC-002} (MEDIUM): ``provably\_secure\_claim'' + ``heuristic\_only'' $\Rightarrow$ QUARANTINE
    \item \textbf{HC-LEGAL-003} (MEDIUM): ``gdpr\_compliant\_claim'' + ``no\_scc'' $\Rightarrow$ QUARANTINE
\end{itemize}

\subsection{Layer 3: Integration (Binary Verdict Gate)}

The integration layer is a pure function:

\begin{algorithm}
\caption{Binary Verdict Gate}
\begin{algorithmic}[1]
\STATE \textbf{Input:} Adjudication state $A$, QI-INT score $S_c$, Configuration $\Gamma$
\STATE \textbf{Output:} Verdict $V \in \{\texttt{SHIP}, \texttt{NO\_SHIP}\}$
\STATE
\IF{$A.\text{kill\_triggered} = \text{true}$}
    \RETURN \texttt{NO\_SHIP} (KILL\_SWITCH)
\ELSIF{$A.\text{rule\_kill} = \text{true}$}
    \RETURN \texttt{NO\_SHIP} (CONTRADICTION\_HIGH\_LEGAL)
\ELSIF{$|A.\text{obligations\_open}| > 0$}
    \RETURN \texttt{NO\_SHIP} (OBLIGATIONS\_BLOCKING)
\ELSIF{$|A.\text{contradictions}| > 0$}
    \RETURN \texttt{NO\_SHIP} (CONTRADICTION\_PRESENT)
\ELSIF{$S_c \geq \Gamma.\tau_{\text{accept}}$}
    \RETURN \texttt{SHIP} (SCORE\_PASS)
\ELSE
    \RETURN \texttt{NO\_SHIP} (SCORE\_FAIL)
\ENDIF
\end{algorithmic}
\end{algorithm}

% ==============================================================================
% 5. QI-INT v2: QUANTUM INTERFERENCE INTEGRATION
% ==============================================================================
\section{QI-INT v2: Complex Amplitude Scoring}

While signals do not vote, they can be mapped to diagnostic amplitudes for score visualization:

\subsection{Mathematical Formulation}

For a claim $C$ at tier $t$ receiving signals from teams $\mathcal{T}$:

\[
a_{c,k} = r_t \cdot e^{i\theta_k}
\]

where:
\begin{itemize}
    \item $r_t \in [0,1]$ is tier magnitude (Tier I = 1.0, Tier II = 0.7, Tier III = 0.4)
    \item $\theta_k$ is signal phase (APPROVE = 0, OBJECT = $\pi/2$, REJECT = $\pi$)
\end{itemize}

Weighted sum:
\[
A_c = \sum_{k \in \mathcal{T}} w_k \cdot a_{c,k}
\]

Score:
\[
S_c = |A_c|^2 = (\Re(A_c))^2 + (\Im(A_c))^2
\]

\subsection{Thresholds}

\begin{itemize}
    \item $\tau_{\text{accept}} = 0.75$: Minimum score for \texttt{SHIP} (if no blockers)
    \item $\tau_{\text{quarantine}} = 0.40$: Reference threshold (diagnostic only)
\end{itemize}

\textbf{Critical property}: QI-INT score is used \textit{only if no kill-switches or obligations are active}. It is a diagnostic tool, not a decision input when hard constraints apply.

% ==============================================================================
% 6. DETERMINISTIC HASHING & REPLAY
% ==============================================================================
\section{Replay Determinism}

\subsection{Canonicalization Protocol}

To ensure deterministic hashing:

\begin{enumerate}
    \item \textbf{Sort signals by team} (alphabetical)
    \item \textbf{Exclude dynamic fields}:
    \begin{itemize}
        \item \texttt{run\_id}, \texttt{timestamp\_start}, \texttt{timestamp\_end}
        \item \texttt{signals[].timestamp}
        \item \texttt{event\_log[].t}
    \end{itemize}
    \item \textbf{Sort JSON keys} (alphabetical)
    \item \textbf{Use fixed-precision arithmetic} (round to 4 decimals)
\end{enumerate}

\subsection{RunManifest Structure}

Every decision produces an immutable \texttt{RunManifest}:

\begin{lstlisting}[caption=RunManifest schema]
{
  "run_id": "uuid::...",
  "code_version": "git::<COMMIT_HASH>",
  "claim": {...},
  "evidence": [...],
  "signals": [...],
  "derived": {
    "kill_switch_triggered": bool,
    "rule_kill_triggered": bool,
    "obligations_open": [...],
    "contradictions": [...],
    "qi_int": {"A_c": {...}, "S_c": float}
  },
  "decision": {
    "final": "SHIP" | "NO_SHIP",
    "ship_permitted": bool,
    "reason_codes": [...]
  },
  "hashes": {
    "inputs_hash": "sha256:...",
    "outputs_hash": "sha256:..."
  }
}
\end{lstlisting}

\begin{theorem}[Replay Equivalence]
For any two executions $E_1, E_2$ with identical canonicalized inputs:
\[
\text{Hash}(\text{inputs}_{E_1}) = \text{Hash}(\text{inputs}_{E_2}) \implies \text{Hash}(\text{outputs}_{E_1}) = \text{Hash}(\text{outputs}_{E_2})
\]
\end{theorem}

% ==============================================================================
% 7. TEST VAULT VALIDATION
% ==============================================================================
\section{Test Vault: Adversarial Validation}

We validate ORACLE SUPERTEAM using 10 adversarial scenarios:

\begin{table}[h]
\centering
\caption{Test Vault Scenarios and Results}
\begin{tabular}{@{}lllc@{}}
\toprule
\textbf{ID} & \textbf{Scenario} & \textbf{Verdict} & \textbf{Pass} \\ \midrule
S-01 & Impossible ROI (500\% in 2 weeks) & NO\_SHIP & \checkmark \\
S-02 & Privacy contradiction + KILL vote & NO\_SHIP (KILL) & \checkmark \\
S-03 & Fake proof (provably secure claim) & NO\_SHIP & \checkmark \\
S-04 & Facial recognition without consent & NO\_SHIP (RULE-KILL) & \checkmark \\
S-05 & Medical claim, no peer review & NO\_SHIP & \checkmark \\
S-06 & Quantum-proof marketing claim & NO\_SHIP & \checkmark \\
S-07 & GDPR claim without SCCs & NO\_SHIP & \checkmark \\
S-08 & \textbf{Replay determinism test} & \textbf{SHIP} & \checkmark \\
S-09 & Kill-switch dominance & NO\_SHIP (KILL) & \checkmark \\
S-10 & Multiple CONDITIONAL → deadlock & NO\_SHIP & \checkmark \\
\bottomrule
\end{tabular}
\end{table}

\textbf{Key result}: 100\% pass rate, with S-08 demonstrating hash-identical outputs for shuffled signal ordering.

% ==============================================================================
% 8. THREAT MODEL
% ==============================================================================
\section{Threat Model}

ORACLE SUPERTEAM is designed to resist:

\subsection{Sovereignty Leaks}

\textbf{Threat}: Agent teams gain decision authority through persuasive language or voting.

\textbf{Mitigation}: Strict signal-only production layer. No agent output can directly influence verdict. All decision logic is hardcoded in adjudication and integration layers.

\subsection{Consensus Gaming}

\textbf{Threat}: Strategic voting or rhetorical inflation to manipulate scores.

\textbf{Mitigation}: QI-INT score is diagnostic only when kill-switches or obligations are active. Lexicographic veto prevents gaming.

\subsection{Replay Instability}

\textbf{Threat}: Non-determinism allows different verdicts for identical inputs.

\textbf{Mitigation}: Deterministic canonicalization and cryptographic hashing. CI enforces replay equivalence on every commit.

\subsection{Evidence Fabrication}

\textbf{Threat}: Fake receipts or unverifiable evidence.

\textbf{Mitigation}: \texttt{NO\_RECEIPT = NO\_SHIP} + hash verification. Future extensions include zero-knowledge proofs for privacy-preserving evidence.

% ==============================================================================
% 9. RELATED WORK
% ==============================================================================
\section{Related Work}

\subsection{Multi-Agent Systems}

\textbf{AutoGen} \cite{autogen} and \textbf{ChatDev} \cite{chatdev}: Conversational agent frameworks with emergent consensus. Lack determinism and hard constraints.

\textbf{MetaGPT} \cite{metagpt}: Software engineering agents with role specialization. Non-deterministic outputs, no veto mechanisms.

\subsection{Formal Verification}

\textbf{Coq}, \textbf{Isabelle/HOL}: Proof assistants with deterministic verification. Non-agentic, require formal specifications.

\subsection{Decision Theory}

\textbf{Voting theory}: Mechanism design for aggregation. Our approach rejects voting in favor of jurisdictional veto.

\textbf{Blockchain governance}: On-chain voting systems. ORACLE SUPERTEAM enforces off-chain determinism without consensus overhead.

% ==============================================================================
% 10. FUTURE WORK
% ==============================================================================
\section{Future Work}

\begin{enumerate}
    \item \textbf{Zero-knowledge receipts}: Privacy-preserving evidence verification
    \item \textbf{Federated oracles}: Multi-organization consensus networks
    \item \textbf{Meta-agent tuning}: Automated parameter optimization with constitutional constraints
    \item \textbf{Human attestation UI}: Web interface for evidence submission and review
    \item \textbf{Parallel Superteams}: ChatDev-style agent teams for upstream production
\end{enumerate}

% ==============================================================================
% 11. CONCLUSION
% ==============================================================================
\section{Conclusion}

ORACLE SUPERTEAM demonstrates that multi-agent systems can be reframed from conversations into \textit{institutions}—computational tribunals with hard constraints, deterministic rules, and cryptographic auditability.

By enforcing non-sovereign production, lexicographic veto, and replay determinism, we achieve governance suitable for high-stakes domains: legal compliance, AI safety gates, research validation, and organizational decision-making.

The framework is production-ready, open-source, and constitutionally constrained.

\textbf{ORACLE SUPERTEAM is not a conversation. It is an institution.}

% ==============================================================================
% ACKNOWLEDGMENTS
% ==============================================================================
\section*{Acknowledgments}

This work was developed independently as part of JMT Consulting's research into epistemic systems and governance automation. The author thanks the open-source community for tools enabling deterministic, auditable software.

% ==============================================================================
% REFERENCES
% ==============================================================================
\begin{thebibliography}{9}

\bibitem{autogen}
Wu, Q., et al. (2023).
\textit{AutoGen: Enabling Next-Gen LLM Applications via Multi-Agent Conversation}.
arXiv preprint arXiv:2308.08155.

\bibitem{chatdev}
Qian, C., et al. (2023).
\textit{ChatDev: Communicative Agents for Software Development}.
arXiv preprint arXiv:2307.07924.

\bibitem{metagpt}
Hong, S., et al. (2023).
\textit{MetaGPT: Meta Programming for Multi-Agent Collaborative Framework}.
arXiv preprint arXiv:2308.00352.

\bibitem{blockchain-gov}
Buterin, V. (2014).
\textit{DAOs, DACs, DAs and More: An Incomplete Terminology Guide}.
Ethereum Foundation Blog.

\bibitem{voting-theory}
Arrow, K. J. (1950).
\textit{A Difficulty in the Concept of Social Welfare}.
Journal of Political Economy, 58(4), 328-346.

\bibitem{zkp}
Goldwasser, S., Micali, S., \& Rackoff, C. (1989).
\textit{The Knowledge Complexity of Interactive Proof Systems}.
SIAM Journal on Computing, 18(1), 186-208.

\end{thebibliography}

% ==============================================================================
% APPENDIX
% ==============================================================================
\appendix

\section{Implementation Details}

The complete implementation is available at:

\begin{center}
\url{https://github.com/yourusername/oracle-superteam}
\end{center}

\textbf{Repository structure}:
\begin{itemize}
    \item \texttt{oracle/}: Core engine (11 Python modules, ~500 LOC)
    \item \texttt{test\_vault/}: 10 adversarial test scenarios
    \item \texttt{ci/}: Continuous integration test runner
    \item \texttt{viz/}: Graphviz governance city map
    \item \texttt{paper/}: This document (LaTeX source)
\end{itemize}

\textbf{Dependencies}: None. Pure Python 3.8+ standard library.

\section{Team Weight Configuration}

\begin{table}[h]
\centering
\caption{Team Weights ($w_k$) and Authority}
\begin{tabular}{@{}lcc@{}}
\toprule
\textbf{Team} & \textbf{Weight} & \textbf{Kill Authority} \\ \midrule
Legal Office & 1.00 & \checkmark \\
Security Sector & 1.00 & \checkmark \\
Engineering Wing & 0.85 & $\times$ \\
Data Validation Office & 0.80 & $\times$ \\
COO HQ & 0.75 & $\times$ \\
UX/Impact Bureau & 0.70 & $\times$ \\
Strategy HQ & 0.65 & $\times$ \\
\bottomrule
\end{tabular}
\end{table}

\section{Graphviz City Map}

The governance topology is visualized as ORACLE TOWN:

\begin{figure}[h]
\centering
\includegraphics[width=0.9\textwidth]{oracle_town.png}
\caption{ORACLE TOWN: Governance city map showing signal flows, kill-switch paths (red), and obligation vault (orange).}
\label{fig:city-map}
\end{figure}

\end{document}
